\documentclass[12pt]{article}

% LuaLaTeX: fontspec + polyglossia для корректной работы с кириллицей
\usepackage{fontspec}
\defaultfontfeatures{Ligatures=TeX,Scale=MatchLowercase}
\setmainfont{Alice-Regular}[Path=font/,Extension=.ttf]
\newfontfamily\cyrillicfont{Alice-Regular}[Path=font/,Extension=.ttf]
\usepackage{polyglossia}
\setmainlanguage{russian}

\usepackage[style=numeric, backend=biber]{biblatex} % Обязательный пакет для корректной сборки
\usepackage{graphicx}        						% Графика
\usepackage{microtype}       						% Микротипографика
\usepackage{amssymb}        						% Математические символы
\usepackage{amsmath}        						% Математика
\usepackage{longtable}      						% Таблицы с разрывом страниц
\usepackage{tikz}           						% Рисование
\usepackage{circuitikz}     						% Электрические схемы
\usepackage{fix-cm}         						% Чтобы работал большой размер шрифта
\usepackage{anyfontsize}    						% Для любого размера шрифта
\usepackage{fancyhdr}        						% Колонтитулы

% Дополнительные размеры для основных шрифтов
\newfontfamily\titlelolita[Scale=1.4]{lolita}[Path=font/,Extension=.ttf]
\newfontfamily\regularlolita[Scale=1.2]{lolita}[Path=font/,Extension=.ttf]
\newfontfamily\titlealice[Scale=1.4]{Alice-Regular}[Path=font/,Extension=.ttf]
\newfontfamily\regularalice[Scale=1.2]{Alice-Regular}[Path=font/,Extension=.ttf]

\newfontfamily\hywenhei[Scale=0.8]{HYWenHei}[Path=font/,Extension=.ttf]

\usepackage[a4paper,left=20mm,right=20mm,top=20mm,bottom=20mm]{geometry}

\title{\titlelolita Дифференциальное Исчисление}
\author{Андреев Иван ЭН-$25$}
\date{}

\begin{document}

\maketitle

\pagebreak

\section{\hywenhei Дифференцирование}

\subsection{\hywenhei Определение дифференциала}

Пусть функция $y = f(x)$ определена в некоторой 
окрестности точки $x$. Приращение аргумента $\Delta x$ 
приводит к приращению функции 
$\Delta y = f(x + \Delta x) - f(x)$. 
Если существует такое выражение вида
\[
\Delta y = f'(x)\Delta x + o(\Delta x),
\]
где $f'(x)$ — некоторая величина, зависящая от $x$, а 
$o(\Delta x)$ — бесконечно малая величина по сравнению с 
$\Delta x$ при $\Delta x \to 0$, то функция $f(x)$ 
называется дифференцируемой в точке $x$, а величина 
$f'(x)\Delta x$ называется дифференциалом функции $f(x)$ 
в этой точке и обозначается $dy$. Таким образом,
\[dy = f'(x)\Delta x,\]
а производная функции $f(x)$ в точке $x$ определяется как
\[f'(x) = \frac{dy}{dx} = \lim_{\Delta x \to 0}\frac{f(x + \Delta x) - f(x)}{\Delta x}.\]

\subsection{\hywenhei Дифференциалы стандартных функций}

\renewcommand{\arraystretch}{3}
\setlength{\tabcolsep}{70pt}
\begin{longtable}{|c|c|}
\hline
Функция $y = f(x)$ & Дифференциал $dy$ \\
\hline
$\displaystyle y = C$ & $\displaystyle dy = 0$ \\
\hline
$\displaystyle y = x^n$ & $\displaystyle dy = nx^{n-1}dx$ \\
\hline
$\displaystyle y = e^x$ & $\displaystyle dy = e^xdx$ \\
\hline
$\displaystyle y = a^x$ & $\displaystyle dy = a^x \ln{a} \, dx$ \\
\hline
$\displaystyle y = \ln{x}$ & $\displaystyle dy = \dfrac{dx}{x}$ \\
\hline
$\displaystyle y = \log_a{x}$ & $\displaystyle dy = \dfrac{dx}{x \ln{a}}$ \\
\hline
$\displaystyle y = \sin{x}$ & $\displaystyle dy = \cos{x} \, dx$ \\
\hline
$\displaystyle y = \cos{x}$ & $\displaystyle dy = - \sin{x} \, dx$ \\
\hline
$\displaystyle y = \tan{x}$ & $\displaystyle dy = \dfrac{dx}{\cos^2{x}}$ \\
\hline
$\displaystyle y = \cot{x}$ & $\displaystyle dy = - \dfrac{dx}{\sin^2{x}}$ \\
\hline
$\displaystyle y = \arcsin{x}$ & $\displaystyle dy = \dfrac{dx}{\sqrt{1 - x^2}}$ \\
\hline
$\displaystyle y = \arccos{x}$ & $\displaystyle dy = - \dfrac{dx}{\sqrt{1 - x^2}}$ \\
\hline
$\displaystyle y = \arctan{x}$ & $\displaystyle dy = \dfrac{dx}{1 + x^2}$ \\
\hline
$\displaystyle y = \operatorname{arccot}{x}$ & $\displaystyle dy = - \dfrac{dx}{1 + x^2}$ \\
\hline
$\displaystyle y = \operatorname{sh}{x}$ & $\displaystyle dy = \operatorname{ch}{x} \, dx$ \\
\hline
$\displaystyle y = \operatorname{ch}{x}$ & $\displaystyle dy = \operatorname{sh}{x} \, dx$ \\
\hline
$\displaystyle y = \operatorname{th}{x}$ & $\displaystyle dy = \dfrac{dx}{\operatorname{ch}^2{x}}$ \\
\hline
$\displaystyle y = \operatorname{cth}{x}$ & $\displaystyle dy = - \dfrac{dx}{\operatorname{sh}^2{x}}$ \\
\hline
$\displaystyle y = \operatorname{arcsh}{x}$ & $\displaystyle dy = \dfrac{dx}{\sqrt{1 + x^2}}$ \\
\hline
$\displaystyle y = \operatorname{arcch}{x}$ & $\displaystyle dy = \dfrac{dx}{\sqrt{x^2 - 1}}$ \\
\hline
$\displaystyle y = \operatorname{arcth}{x}$ & $\displaystyle dy = \dfrac{dx}{1 - x^2}$ \\
\hline
$\displaystyle y = \operatorname{arcth}{x}$ & $\displaystyle dy = - \dfrac{dx}{1 - x^2}$ \\
\hline
\end{longtable}

\begin{enumerate}
	\item {\hywenhei Найти дифференциал функции $y = C$} \\
	По определению производной имеем
	\[
		f'(x) = \lim_{\Delta x \to 0}\frac{C - C}{\Delta x} = 0 \implies dy = 0 \cdot dx = 0
	\]
	\item {\hywenhei Найти дифференциал функции $y = x^a$} \\
	Расспишем Бином Ньютона
	\[
		\displaystyle (a + b)^a = \sum_{k = 0}^{n = a}\Big[ a^{n-k} b^k C_k^n \Big]\,,\text{ где } C_k^n = \dfrac{n!}{(n-k)!k!}
	\]
	Тогда по определению производной имеем
	\[
		\begin{aligned}
			\displaystyle f'(x) &= \lim_{\Delta x \to 0}\dfrac{(x + \Delta x)^a - x^a}{\Delta x} = \lim_{\Delta x \to 0}\dfrac{\displaystyle x^a + \sum_{k = 1}^{n = a} \Big[ x^{n-k} (\Delta x)^k C_k^n \Big] - x^a}{\Delta x} = \\
								&= \lim_{\Delta x \to 0}\dfrac{\displaystyle \sum_{k = 1}^{n = a} \Big[ x^{n-k} (\Delta x)^k C_k^n \Big]}{\Delta x} = \lim_{\Delta x \to 0}\sum_{k = 1}^{n = a} \Big[ x^{n-k} (\Delta x)^{k-1} \frac{n!}{(n-k)!k!} \Big] = \\
								&= \lim_{\Delta x \to 0}\left(ax^{a-1}+\frac{a}{2}x^{a-2}\Delta x+\dots+ax\Delta x^{a-2}+\Delta x^{a-1}\right)=\lim_{\Delta x \to 0}\Big[ax(x+\Delta x)^{a-2}\Big]=ax^{a-1}
		\end{aligned}
	\]
	Возможно также пропустить предпоследний шаг, так как все члены, кроме первого, при $\Delta x \to 0$ стремятся к нулю. Тогда дифференциал будет равен
	\[
		dy = ax^{a-1}dx
	\]
	\item {\hywenhei Найти дифференциал функции $y = e^x$} \\
	По определению производной имеем
	\[
		f'(x) = \lim_{\Delta x \to 0}\frac{e^{x + \Delta x} - e^x}{\Delta x} = \lim_{\Delta x \to 0}\frac{e^x(e^{\Delta x} - 1)}{\Delta x} = e^x \lim_{\Delta x \to 0}\frac{e^{\Delta x} - 1}{\Delta x} = e^x \lim_{\Delta x \to 0}\frac{\Delta x}{\Delta x} = e^x
	\]
	Тогда дифференциал будет равен
	\[
		dy = e^x dx
	\]
\item {\hywenhei Найти дифференциал функции $y = a^x$} \\
	По определению производной имеем
	\[
		\begin{aligned}
			f'(x) &= \lim_{\Delta x \to 0}\frac{a^{x + \Delta x} - a^x}{\Delta x} = \lim_{\Delta x \to 0}\frac{a^x(a^{\Delta x} - 1)}{\Delta x} = a^x \lim_{\Delta x \to 0}\frac{a^{\Delta x} - 1}{\Delta x} = \\
				  &= a^x \lim_{\Delta x \to 0}\frac{e^{\Delta x \ln{a}} - 1}{\Delta x} = a^x \lim_{\Delta x \to 0}\frac{\Delta x \ln{a}}{\Delta x} = a^x \ln{a}
		\end{aligned}
	\]
	Тогда дифференциал будет равен
	\[
		dy = a^x \ln{a} \, dx
	\]
\item {\hywenhei Найти дифференциал функции $y = \ln{x}$} \\
	По определению производной имеем
	\[
		f'(x) = \lim_{\Delta x \to 0}\frac{\ln{(x + \Delta x)} - \ln{x}}{\Delta x} = \lim_{\Delta x \to 0}\frac{\ln{\frac{x + \Delta x}{x}}}{\Delta x} = \lim_{\Delta x \to 0}\frac{\ln{(1 + \frac{\Delta x}{x})}}{\Delta x}= \lim_{\Delta x \to 0}\frac{\frac{\Delta x}{x}}{\Delta x} = \frac{1}{x}
	\]
	Тогда дифференциал будет равен
	\[
		dy = \dfrac{dx}{x}
	\]
\item {\hywenhei Найти дифференциал функции $y = \log_a{x}$} \\
	По определению производной имеем
	\[
		\begin{aligned}
			f'(x) &= \lim_{\Delta x \to 0}\frac{\log_a{(x + \Delta x)} - \log_a{x}}{\Delta x} = \lim_{\Delta x \to 0}\frac{\log_a{\frac{x + \Delta x}{x}}}{\Delta x} = \lim_{\Delta x \to 0}\frac{\log_a{(1 + \frac{\Delta x}{x})}}{\Delta x}= \\
				  &= \lim_{\Delta x \to 0}\frac{\frac{\ln{(1 + \frac{\Delta x}{x})}}{\ln{a}}}{\Delta x} = \frac{1}{\ln{a}} \lim_{\Delta x \to 0}\frac{\ln{(1 + \frac{\Delta x}{x})}}{\Delta x} = \frac{1}{\ln{a}} \lim_{\Delta x \to 0}\frac{\frac{\Delta x}{x}}{\Delta x} = \frac{1}{x \ln{a}}
		\end{aligned}
	\]
	Тогда дифференциал будет равен
	\[
		dy = \dfrac{dx}{x \ln{a}}
	\]
\item {\hywenhei Найти дифференциал функции $y = \sin{x}$} \\
	По определению производной имеем
	\[
	\begin{aligned}
		f'(x) &= \lim_{\Delta x \to 0}\frac{\sin{(x + \Delta x)} - \sin{x}}{\Delta x} = \lim_{\Delta x \to 0}\frac{\sin{x}\cos{\Delta x} + \cos{x}\sin{\Delta x} - \sin{x}}{\Delta x} = \\
			  &= \lim_{\Delta x \to 0}\frac{\sin{x}(\cos{\Delta x} - 1) + \cos{x}\sin{\Delta x}}{\Delta x} = \sin{x} \lim_{\Delta x \to 0}\frac{\cos{\Delta x} - 1}{\Delta x} + \cos{x} \lim_{\Delta x \to 0}\frac{\sin{\Delta x}}{\Delta x} = \\
			  &= \sin{x} \lim_{\Delta x \to 0}\frac{\sin^2(\frac{\Delta x}{2})}{\frac{\Delta x}{2}} + \cos{x} = \sin{x} \cdot 0 + \cos{x} = \cos{x}
	\end{aligned}
	\]
	Тогда дифференциал будет равен
	\[
		dy = \cos{x} \, dx
	\]
\item {\hywenhei Найти дифференциал функции $y = \cos{x}$} \\
	По определению производной имеем
	\[
	\begin{aligned}
		f'(x) &= \lim_{\Delta x \to 0}\frac{\cos{(x + \Delta x)} - \cos{x}}{\Delta x} = \lim_{\Delta x \to 0}\frac{\cos{x}\cos{\Delta x} - \sin{x}\sin{\Delta x} - \cos{x}}{\Delta x} = \\
			  &= \lim_{\Delta x \to 0}\frac{\cos{x}(\cos{\Delta x} - 1) - \sin{x}\sin{\Delta x}}{\Delta x} = \cos{x} \lim_{\Delta x \to 0}\frac{\cos{\Delta x} - 1}{\Delta x} - \sin{x} \lim_{\Delta x \to 0}\frac{\sin{\Delta x}}{\Delta x} = \\
			  &= \cos{x} \lim_{\Delta x \to 0}\frac{\sin^2(\frac{\Delta x}{2})}{\frac{\Delta x}{2}} - \sin{x} = \cos{x} \cdot 0 - \sin{x} = -\sin{x}
	\end{aligned}
	\]
	Тогда дифференциал будет равен
	\[
		dy = - \sin{x} \, dx
	\]
\item {\hywenhei Найти дифференциал функции $y = \tan{x}$} \\
	По определению производной имеем
	\[
	\begin{aligned}
		f'(x) &= \lim_{\Delta x \to 0}\frac{\tan{(x + \Delta x)} - \tan{x}}{\Delta x} = \lim_{\Delta x \to 0}\frac{\frac{\sin{(x + \Delta x)}}{\cos{(x + \Delta x)}} - \frac{\sin{x}}{\cos{x}}}{\Delta x} = \\
			  &= \lim_{\Delta x \to 0}\frac{\sin{(x + \Delta x)}\cos{x} - \sin{x}\cos{(x + \Delta x)}}{\Delta x \cdot \cos{(x + \Delta x)}\cos{x}} = \\
			  &= \lim_{\Delta x \to 0}\frac{\sin x\cos\Delta x\cos x + \cos^2 x\sin\Delta x - \cos x\cos\Delta x\sin x + \sin^2 x\sin \Delta x}{(\cos x\cos\Delta x - \sin x\sin\Delta x)\Delta x\cos x} = \\
			  &= \lim_{\Delta x \to 0}\frac{\cos^2 x\sin\Delta x + \sin^2 x\sin \Delta x}{(\cos x\cos\Delta x - \sin x\sin\Delta x)\Delta x\cos x} = \\
			  &= \lim_{\Delta x \to 0}\frac{\sin\Delta x}{\Delta x} \cdot \lim_{\Delta x \to 0}\frac{\cos^2 x + \sin^2 x}{(\cos x\cos\Delta x - \sin x\sin\Delta x)\cos x} = \\
			  &= 1 \cdot \frac{1}{(1\cdot\cos x - 0\cdot\sin x)\cos x} = \frac{1}{\cos^2 x}
	\end{aligned}
	\]
	Тогда дифференциал будет равен
	\[
		dy = \dfrac{dx}{\cos^2{x}}
	\]
\item {\hywenhei Найти дифференциал функции $y = \cot{x}$} \\
	По определению производной имеем
	\[
	\begin{aligned}
		f'(x) &= \lim_{\Delta x \to 0}\frac{\cot{(x + \Delta x)} - \cot{x}}{\Delta x} = \lim_{\Delta x \to 0}\frac{\frac{\cos{(x + \Delta x)}}{\sin{(x + \Delta x)}} - \frac{\cos{x}}{\sin{x}}}{\Delta x} = \\
			  &= \lim_{\Delta x \to 0}\frac{\cos{(x + \Delta x)}\sin{x} - \cos{x}\sin{(x + \Delta x)}}{\Delta x \cdot \sin{(x + \Delta x)}\sin{x}} = \\
			  &= \lim_{\Delta x \to 0}\frac{\cos x\cos\Delta x\sin x - \sin^2 x\sin\Delta x - \cos x\cos\Delta x\sin x - \cos^2 x\sin\Delta x}{(\sin x\cos\Delta x + \cos x\sin\Delta x)\Delta x\sin x} = \\
			  &= \lim_{\Delta x \to 0}\frac{-\sin^2 x\sin\Delta x - \cos^2 x\sin\Delta x}{(\sin x\cos\Delta x + \cos x\sin\Delta x)\Delta x\sin x} = \\
			  &= \lim_{\Delta x \to 0}\frac{-\sin\Delta x}{\Delta x} \cdot \lim_{\Delta x \to 0}\frac{\sin^2 x + \cos^2 x}{(\sin x\cos\Delta x + \cos x\sin\Delta x)\sin x} = \\
			  &= -1 \cdot \frac{1}{(1\cdot\sin x + 0\cdot\cos x)\sin x} = -\frac{1}{\sin^2 x}
	\end{aligned}
	\]
	Тогда дифференциал будет равен
	\[
		dy = - \dfrac{dx}{\sin^2{x}}
	\]
\end{enumerate}

\subsection{\hywenhei Дифференциальные уравнения}

TODO: продемеонстрировать виды уравнений и методы их решения

\section{\hywenhei Интегрирование}

\subsection{\hywenhei Определение интеграла Римана}

Пусть функция $f(x)$ непрерывна на отрезке $[\,a,\, b\,]$ . Разобьём этот отрезок на $n$ частей 
точками $x_0,\, x_1,\, \ldots,\, x_n$ , где $a = x_0 < x_1 < \ldots < x_n = b$ . Обозначим через 
$\Delta x_i = x_{i + 1} - x_i$ длину $i$-го отрезка разбиения и через
$\xi_i$ произвольную точку этого отрезка, то есть $x_i \leq \xi_i \leq x_{i + 1}$ . Построим 
сумму 
\[S_n = \sum_{i=0}^{n-1} f(\xi_i) \Delta x_i.\]
Если при стремлении максимальной длины отрезков разбиения
\[\lambda = \max_{0 \leq i \leq n} \Delta x_i \to 0\]
существует предел суммы $S_n$ , не зависящий от выбора точек $\xi_i$ , то этот предел называется 
определённым интегралом функции $f(x)$ на отрезке $[\,a,\, b\,]$ и обозначается
\[S = \lim_{\lambda \to 0} S_n = \int\limits_a^b f(x) \, dx.\]

\noindent {\hywenhei Теорема о среднем.} Если функция $f(x)$ непрерывна на отрезке 
$[\,a,\, b\,]$ , то существует такая точка $\xi \in [\,a,\, b\,]$ , что
\[\int\limits_a^b f(x) \, dx = f(\xi)(b - a).\]

\noindent {\hywenhei Основная Теорема Анализа.} Если функция $f(x)$ непрерывна 
на отрезке $[\,a,\, b\,]$ , то верно
\[f(x) = \left(\int\limits_{a}^{x}f(t)\,dt\right)'.\]

\subsection{\hywenhei Интегралы стандартных функций}

\renewcommand{\arraystretch}{3}
\setlength{\tabcolsep}{50pt}
\begin{longtable}{|c|c|}
	\hline
Функция $f(x)$ 
& 
Неопределённый интеграл $\displaystyle \int f(x) \, dx$ \\
\hline
$\displaystyle f(x) = \tilde{C}$ 
&
$\displaystyle \int f(x) \, dx = \tilde{C}x + C$ \\
\hline
$\displaystyle f(x) = x^a$ 
&
$\displaystyle \int f(x) \, dx = \frac{x^{a+1}}{a+1} + C$ \\
\hline
$\displaystyle f(x) = e^x$
&
$\displaystyle \int f(x) \, dx = e^x + C$ \\
\hline
$\displaystyle f(x) = a^x$
&
$\displaystyle \int f(x) \, dx = \frac{a^x}{\ln{a}} + C$ \\
\hline
\end{longtable}

\begin{enumerate}
\item {\hywenhei Найти неопределённый интеграл функции $f(x) = \tilde{C}$} \\
TODO: добавить решение
\end{enumerate}

\end{document}